
\chapter{Static App Analysis}
\label{chap:appAnalysis}

This chapter presents a static analysis of the Meta View app, the companion app for the Ray-Ban Meta glasses. The analysis focuses on examining important aspects of the app that would aid in understanding its role in managing user data and its privacy implications. These aspects include the app's requested permissions, the integration of third-party libraries, and the potential transmission and recipient of user information.

The app serves as the primary interface between the glasses and Meta's backend services. Since it manages core functionalities, such as pairing, media import, and notifications, the permissions requested by the app and its libraries can significantly impact user privacy. Therefore, examining the permissions and dependencies is important to understand the potential risks related to data collection, user tracking, or unauthorized access. The static analysis was performed using MobSF to inspect the app's APK, providing details about permissions, native libraries, manifest components, and other relevant information.


\section{Permissions Analysis}
This section covers the permissions identified in the app as well as the privacy and security risks when granted. Android applications must declare which system resources they require through permissions. The Meta View app requests several permissions that are considered dangerous by the Android operating system. Dangerous permissions grant access to sensitive user data or critical features, making their presence noteworthy when evaluating potential security and privacy risks. The following dives into these permissions and their potential implications based on the permission types.

\subsection*{Location}
The app requests multiple location-related permissions, including background, coarse (network-based), and fine (GPS). While these location permissions are necessary for functionalities such as pairing with the Ray-Ban Meta glasses or providing location-based services, they also raise some concerns. Background location permission allows the app to track the user's whereabouts even whilst the app is not actively in use, which creates the possibility of constant user tracking. Coarse location access uses network data, for example, cell towers, to determine an approximate location of the user, whilst fine location uses GPS to determine a precise location, thereby allowing for higher accuracy in the mapping of user movements. These permissions could allow third parties to track the precise movement of a user by continuously collecting GPS data, essentially like physical surveillance. Furthermore, when combined with other permissions, such as network access, these location permissions could lead to detailed profiling of user activities and movements.


\subsection*{Bluetooth}
Permissions to access Bluetooth-related functionalities include scanning and connection, which allows scanning for, discovering, and pairing with nearby devices. While necessary for pairing with the Ray-Ban glasses, these permissions present a risk where malicious devices within range could attempt to establish unauthorized pairing and access to the glasses. Furthermore, if the connections are not properly secured, data exchanged via Bluetooth could be intercepted or tampered with by a third party.

\subsection*{Account and Device Connectivity}
In order to access user accounts and manage the device's network connections, the app requires permissions such as GET\_ACCOUNTS and access to nearby Wi-Fi devices. These permissions support importing images from the glasses to the app. However, if misused, there is a risk that an attacker or the app itself could gain unauthorized access. The Get Accounts permission allows the app to access a list of accounts in the Accounts Service, which could be used for unauthorized data access or data gathering. The permission for the connection to nearby Wi-Fi devices allows the app to advertise and connect to nearby devices, which could expose the user to malicious devices within range attempting to initiate connections over Wi-Fi, potentially leading to eavesdropping or data interceptions by third parties.

\subsection*{Storage and Media}
By requesting read, write, and delete access to data from external storage and media files, the app is able to manage photos and videos captured by the glasses. However, the problem with these permissions is that the storage includes personal data, including photos, videos, or even documents, which, if exploited, could be accessed or modified without user consent. Another concern is that the app could theoretically upload these files to remote servers without explicit user consent, which infringes on data privacy and user autonomy. Furthermore, if this permission is combined with location permissions, photo metadata could be used to track user activity and movements.

\subsection*{Phone and Call}
The app requests permission to read the call log and contact data, allowing the app to read and access the user's contact information and call history. The app also requests permission to read the phone state and identity, allowing the app to access the phone identifiers, including the phone number, serial number, and whether or not the user is on a call. These permissions allow the app to access and collect sensitive information, including details on their communications and social circles. The ability to access phone identifiers raises concerns about user and device fingerprinting across multiple services. If exploited, these permissions could reveal personal information such as the user's communication patterns, call history, relationships, and private contact information, which could be used in targeted attacks. 

\subsection*{Communication and Messaging}
The app requested several communication-related permissions, including receiving, reading, and sending messages. Although necessary for features such as communicating on messaging platforms, these permissions allow access to personal user communications. This access introduces the possibility of unauthorized retrieval of messages or intercepted one-time passwords, which could be used for phishing attacks where the data are used for identity theft or unauthorized messaging. The app also requests permission to post notifications, which allows the app to post notifications. This is a valid request for the posting functionality on social media platforms; however, it could be exploited for deceptive or phishing attempts via fake notifications.

The list of permissions found in the Meta View app is provided in Table \ref{tab:summaryPermissions}. The permissions in Malware are the most widely abused permissions that allow access to highly sensitive data, or critical system features such as storage, location, SMS, and phone state. If abused, it can expose users to threats such as unauthorized data collection or surveillance. Although not as extensive, the permissions in Common are also commonly abused and pose privacy or security risks. The classification between dangerous and normal depends on the level of risk associated with granting these permissions to an app. The Android permission system defines this classification. Permissions classified as dangerous grant access to sensitive information and, therefore, have a higher potential to affect user privacy or security. On the other hand, permissions classified as normal allow access to less sensitive data and are considered low risks to user privacy and security.

In summary, the permissions requested by the Meta View app, as summarized in Table \ref{tab:summaryPermissions}, are necessary for certain app functionality, including managing connectivity with the Ray-Ban glasses and enabling features like photo sharing, notifications, and voice commands. Nevertheless, granting these permissions exposes users to potential privacy violations, such as continuous precise location tracking, data harvesting, or interception of personal communications. 

%%%%%%%%%%%%%%%%%%%%%%%%%%%%%%%%%%%%%%%%%%
\section{Third-Party Usage}
This section examines whether the app includes any third-party libraries, determines whether personal information is being transmitted, and, if so, who the recipient is. The investigation focuses on components of the Meta View app that may involve third-party usage or data transmission, including native libraries, trackers, manifest components, and databases, to identify relevant data flows and privacy implications. 

The libraries and components are also grouped into different party types for the static app analysis. Similar to traffic analysis, first-party refers to those owned by Meta and directly linked to the app. Support parties refer to the components that Meta owns but are external to the app. Third-party are libraries, services, or components that Meta does not own and instead provided by an external organization. Additionally, there is an additional party included here, which is platform-provided parties. They are components provided as part of the Android development framework or operating system (e.g., AndroidX).

\subsection*{Native Libraries}
Native libraries are compiled binaries (.so files), usually written in C or C++, that potentially integrate code from external sources. Thirteen native libraries were identified in the app. A complete and detailed table is provided in \ref{tab:sharedObjectsNativeLibraries}. Among these, only two libraries, libbreakpad and libbreakpad\_cpp\_helper, are third-party libraries. Google developed both libraries for crash reporting. They are typically used for debugging and improving app stability, allowing developers to collect and analyze crash data. The remaining libraries appear to be first-party or platform-provided, determined through the human-readable strings embedded within the APK file. Their functionalities, such as dynamic linking and signal handling, also align closely with tasks typically associated with internal development.

\subsection*{Trackers}
Trackers are mechanisms in an app used to monitor user behavior, usually for analytics or advertising purposes. Therefore, analyzing the trackers integrated into the app can identify potential data collection and transmissions. The Meta View app has the Facebook Flipper tracker, categorized as an analytics tracker. Analytics tracking is designed to collect data on how users interact with the app, such as which features are used the most, how long users stay engaged, or where they encounter problems.

Facebook Flipper is used in the app to transmit user data to Meta or Facebook for analytics purposes. However, although the stated purpose is to help improve the user experience or to find bugs, it raises potential privacy concerns as analytics trackers may collect data that could include potentially sensitive information, especially without sufficient user consent. Although there is no solid evidence that the tracker shares data with external organizations or advertisers, this analytics tracker transmits user engagement information to the internal servers of Meta, which may be combined with other user data for better user profiling.

\subsection*{Manifest Components}
The AndroidManifest.xml file serves as a blueprint for the app, providing metadata that defines the structure, permission, and components of the app. By analyzing the list of components, including services, activities, and receivers, the aim is to identify which components are tied to third-party libraries. The app includes 20 services, receivers, and activities that are exported. These components are accessible not only to the Meta View app but also to other apps on the device. As a result, it is a potential entry point for interactions with third-party libraries. A complete detailed table of these components can be found in \ref{tab:manifest_analysis}. The majority of exported components are first-party or support-party. While these do not involve third-party organizations, they may transmit user data internally within the Meta ecosystem. Only four components were found to be tied to third-party components.

The first component, FirebaseInstanceIdReceiver, is owned by Google and supports push notifications by registering device tokens. It aids in the communication between the app and Firebase Cloud Messaging. This component transmits device tokens and potential notification metadata to Google's Firebase Cloud Messaging servers. The second component, PushLiteGCMJobService, manages lightweight push notifications. Although the service is implemented by Meta, it relies on Google's infrastructure, specifically Google Cloud Messaging, and therefore is classified as third-party. As with the first component, it transmits device tokens and potential notification metadata to Google servers.

For both components, if the app sends sensitive information in the payload, such as account updates regarding password changes, this data could be exposed during transmission or logged on Firebase. Although notifications are usually encrypted, metadata accompanying them could be accessible, including the device model or timestamps of when the notification was sent. Similarly, like the Meta View app, Firebase may also collect additional data for its analytics, which could include the frequency of push notifications, resulting in potential breaches of user privacy.

The third component, RevocationBoundService, is part of a library that allows users to authenticate with their Google accounts in the app. This service is part of Google Play Services, specifically, Google Sign-In API. This library is used for managing sign-in tokens and other related authentication features, meaning that authentication tokens and user account data necessary for authentication and session management are transmitted to Google, most likely including email address and user ID. Authentication tokens and user account details are sensitive and, if intercepted, could allow unauthorized access to Google accounts.

The last third-party component, RedirectUriReceiverActivity, is part of the Spotify Authentication SDK, a library allowing users to authenticate their Spotify account through open authorization (OAuth) tokens. Open authorization tokens and possibly user account details may be transmitted to Spotify servers as part of the authentication process. Similarly with Google, authentication tokens and user account details are sensitive and, if intercepted, could allow unauthorized access to Spotify accounts. 

In summary, the app integrates several third-party libraries, including Firebase (Google), Google Cloud Messaging, Google Sign-In, and Spotify Authentication SDK. These libraries are necessary for features such as push notifications and user authentication. However, they may transmit sensitive data, such as personal information, device tokens, open authorization tokens, user credentials, and account data, to their respective servers. If these transmissions are intercepted by external entities or used by their respective platforms for unintended purposes, there is a risk of unauthorized access.


\subsection*{Firebase Database}
Firebase, developed by Google, provides cloud-based, backend, and analytics services. The services it provides mean that Firebase databases may handle sensitive user data. The app communicates with the Firebase endpoint fb-messaging-b8df1.firebaseio.com. On a positive note, the Firebase Remote Config is disabled in the app, which is a feature that allows dynamic modifications of the app's behavior and appearance directly from the server side without requiring a new version or update of the app to be launched. Firebase Remote Config could expose unauthorized users to sensitive parameters such as API keys. Therefore, with it being disabled, the risk of exposing sensitive user data or security vulnerabilities from unauthorized parameter changes of the app is reduced. However, since the app still communicates with Firebase and the exact data transmitted and stored remain unclear, the possibility that sensitive user-related data could be shared in a third-party environment remains a concern.

\section{Summary}
The Meta View app's analysis reveals that it requests a wide range of permissions, including those classified as dangerous. Although these permissions, such as location access and read or write privileges, are required for the app's functionalities, they expose users to potential risks, including continuous location tracking or unauthorized access to personal media and files.

The Meta View app identified various third-party libraries and potential data flows to third-party organizations, particularly from Google, including Firebase, Breakpad, and Google Sign-In, and Spotify, including Authentication SDK. Although these components provide important functionalities such as crash reporting, push notifications, and user authentication, transmitting sensitive data, such as device tokens and user credentials, to third-party organizations raises privacy concerns. Furthermore, first-party trackers, such as Facebook Flipper and other native libraries and components, are first-party or platform-provided and most likely collect user behavior data and transmit it to Meta analytics servers. Although its existence reflects that the app mainly relies on its internal infrastructure and the foundational tools of the Android development framework, it raises privacy concerns about the scope of internal data sharing within the Meta ecosystem.

The analysis details the importance of clearly communicating these data flows and permission requirements to users for a more informed market and the necessity for robust data protection policies and security measures to minimize the risks. Overall, the findings suggest that the Meta View app functionalities rely on external parties, namely Google and Spotify, and practice data collection and sharing, which potentially risks user privacy.